\begin{center}
\small
\begin{longtable}{|p{2.6cm}|p{2.3cm}|p{3.6cm}|p{3.8cm}|p{3.2cm}|}
\caption{Сравнение существующих аналогов и подходов}\label{tab:analogs_comparison}\\
\hline
\textbf{Подход / средство} & \textbf{Класс} & \textbf{Сильные стороны} & \textbf{Ограничения} & \textbf{Роль в исследовании} \\
\hline
\endfirsthead
\hline
\textbf{Подход / средство} & \textbf{Класс} & \textbf{Сильные стороны} & \textbf{Ограничения} & \textbf{Роль в исследовании} \\
\hline
\endhead
Buy \& Hold & baseline-стратегия & Простота, прозрачность, отражение рыночного бета-эффекта & Не адаптируется к рыночным режимам, не управляет риском активно & Базовый эталон абсолютной доходности \\
\hline
Moving Average Crossover & техническая стратегия & Интерпретируемость, учет тренда, низкая вычислительная сложность & Запаздывание сигнала, чувствительность к параметрам окна, ложные входы на боковом рынке & Сопоставление RL с простым правилом на индикаторах \\
\hline
ARIMA / SARIMA & статистическое прогнозирование & Формальный аппарат, интерпретируемые параметры, удобство для линейной динамики & Плохо описывает нелинейность и смену режимов, требует стационаризации & Аналог прогнозирования ценового ряда \\
\hline
Logistic Regression & классическое МО & Простота, устойчивость, возможность анализа вклада признаков & Ограниченная нелинейность, зависимость от качества признаков & Базовая ML-модель для бинарного прогноза направления \\
\hline
Random Forest & ансамблевое МО & Работа с нелинейными зависимостями, устойчивость к шуму, малая потребность в масштабировании & Снижение интерпретируемости, риск переобучения при слабой валидации & Сравнение с линейной моделью и RL \\
\hline
Gradient Boosting & ансамблевое МО & Высокая гибкость, качественная аппроксимация сложных закономерностей & Требовательность к настройке, чувствительность к качеству данных & Сильный ML-базис в экспериментальном стенде \\
\hline
CatBoost / XGBoost & градиентный бустинг & Высокое качество на табличных данных, эффективная работа с признаками & Выше вычислительные затраты, необходимость тщательного тюнинга & Аналог продвинутого табличного бустинга \\
\hline
LSTM / GRU & нейросетевые модели & Учет последовательной структуры, моделирование временной зависимости & Чувствительность к объему данных, сложность настройки, риск нестабильного обучения & Аналоги глубокого обучения для временных рядов \\
\hline
DQN & RL-алгоритм value-based & Обучение дискретным действиям, прямое связывание состояния и торгового решения & Нестабильность при плохой reward function, чувствительность к коррелированным данным & RL-аналог с дискретным пространством действий \\
\hline
PPO & RL-алгоритм policy-based & Устойчивое обновление политики, хорошая практическая применимость, удобство для continuous optimization & Более высокая стоимость обучения, зависимость от дизайна среды & Основной RL-алгоритм для сравнения риск-скорректированной эффективности \\
\hline
A2C & actor-critic & Более компактная альтернатива PPO, обучает политику и value function совместно & Меньшая устойчивость по сравнению с PPO на шумных рынках & Рассматривается как альтернативный RL-подход в обзоре \\
\hline
FinRL & фреймворк & Готовые среды, примеры RL-трейдинга, акцент на reproducibility & Сильная зависимость от встроенных сценариев, сложность адаптации под конкретный pipeline & Аналог исследовательской платформы \\
\hline
TensorTrade & фреймворк & Модульная архитектура, абстракции обмена, портфеля и вознаграждения & Повышенный порог входа, необходимость глубокой настройки & Аналог конструкторного инструментария \\
\hline
gym-anytrading & библиотека сред & Простота, совместимость с Gym, быстрое прототипирование & Ограниченная гибкость и сравнительно простая рыночная логика & Аналог минималистичной RL-среды \\
\hline
\end{longtable}
\end{center}
